\chapter{צמיחה, הפצה ומנועי צמיחה}
\label{chap:growth}

מוצר מצוין אינו מתפשט מעצמו. פיטר תיל מזהיר שמייסדים טכניים נוטים לזלזל
ב\emph{הפצה} \term{(distribution)} ולהאמין שאיכות המוצר תמכור את עצמה — אמונה
שמכשילה חברות רבות, ודווקא ככל שהמוצר טוב יותר~\cite{thiel2014zero}. בניית מנוע
צמיחה שיטתי חשובה לא פחות מבניית המוצר עצמו.

\section{שלושת מנועי הצמיחה}

אריק ריס מבחין בין שלושה מנועי צמיחה, שכל אחד מהם נמדד אחרת~\cite{ries2011lean}.
המנוע ה\emph{דביק} \term{(sticky)} מסתמך על שימור גבוה: צמיחה מתרחשת כאשר קצב
רכישת הלקוחות עולה על קצב הנטישה. המנוע ה\emph{ויראלי} \term{(viral)} מסתמך על
כך שכל משתמש מביא משתמשים נוספים, ונמדד במקדם הוויראליות
\term{(viral coefficient)}. המנוע ה\emph{בתשלום} \term{(paid)} ממחזר רווח
מהלקוח כדי להשיג לקוחות נוספים, וקיים כל עוד \en{LTV} עולה על \en{CAC} כפי
שנדון בפרק~\ref{chap:economics}. ניסיון להפעיל את שלושת המנועים בו-זמנית מפזר
את המאמץ; עדיף למקד באחד ולמצותו.

\begin{figure}[htbp]
\centering
% English-only node labels inside an LTR (\entikz) context: babel bidi=basic
% otherwise reorders the labels across the picture. Hebrew goes in the caption.
\entikz{%
\begin{tikzpicture}[every node/.style={font=\sffamily,align=center}]
  \node[draw=brandblue, fill=brandblue!8, rounded corners, minimum width=3.2cm,
        minimum height=1.7cm] (sticky) at (0,0)
        {{\bfseries\large\color{brandblue}Sticky}\\[3pt]{\small retention}};
  \node[draw=brandgreen, fill=brandgreen!8, rounded corners, minimum width=3.2cm,
        minimum height=1.7cm] (viral) at (4,0)
        {{\bfseries\large\color{brandgreen}Viral}\\[3pt]{\small viral coefficient}};
  \node[draw=brandred, fill=brandred!8, rounded corners, minimum width=3.2cm,
        minimum height=1.7cm] (paid) at (8,0)
        {{\bfseries\large\color{brandred}Paid}\\[3pt]{\small LTV > CAC}};
\end{tikzpicture}}
\caption{שלושת מנועי הצמיחה — \en{Sticky} (שימור גבוה), \en{Viral} (כל משתמש מביא
משתמשים) ו-\en{Paid} (מימון רכישה מתוך רווח). כל מנוע נמדד במדד-על אחר ומחייב מיקוד נפרד.}
\label{fig:engines}
\end{figure}

\section{התאמת ערוץ למוצר}

לכל מוצר מתאים מספר מצומצם של ערוצי הפצה, ולרוב ערוץ דומיננטי יחיד. גובה מחיר
העסקה הממוצע קובע אילו ערוצים כלכליים עבורו: מוצר בעלות נמוכה אינו יכול לממן
צוות מכירות אנושי, ולכן נשען על צמיחה ``מלמטה למעלה'' \term{(bottom-up)} ועל
מוצר שמוכר את עצמו; מוצר ארגוני יקר מצדיק מכירה ישירה ומחזור מכירה ארוך. טעות
נפוצה היא להעתיק את ערוץ ההפצה של חברה מצליחה אחרת מבלי לבדוק אם מבנה העלויות
והלקוחות דומה כלל.

\section{מצמיחה לאחיזה}

צמיחה ללא אחיזה \term{(retention)} היא דלי מחורר: ככל שמזרימים לקוחות מהר יותר,
כך בולטת הדליפה. לכן המדד הראשון שיש לייצב הוא עקומת השימור — האחוז שנשאר פעיל
לאורך זמן — ורק לאחר שהיא משתטחת ברמה בריאה כדאי ללחוץ על דוושת הגז של הרכישה.
צמיחה מוקדמת מדי, בטרם הושגה אחיזה, מייקרת את הכישלון ומסתירה את הבעיה האמיתית
במקום לפתור אותה.

\begin{takeaway}
בחרו מנוע צמיחה \emph{אחד} שמתאים למבנה העלויות שלכם, ודאו אחיזה לפני שמאיצים,
ואל תעתיקו את ערוץ ההפצה של חברה אחרת — מה שעובד עבורה אינו בהכרח כלכלי עבורכם.
\end{takeaway}
